\section{Introduction}

GPU resource management policies critically influence application performance. These policies include memory placement, migration, and GPU task scheduling. Recent research shows their substantial impact on Unified Virtual Memory (UVM) placement~\cite{wang2024suv,lin2025forest,kim2025dream,park2025helm} and GPU scheduling strategies~\cite{wang2024gcaps,fan2025gpreempt,shen2025xsched}. Under memory oversubscription, for instance, eviction and prefetch policies alone can determine up to 73\% of execution time~\cite{park2025helm,kehne2019etc,ganguly2019interplay}.

Unfortunately, finding an optimal GPU resource management policy is challenging. The optimal memory and scheduling policies vary significantly across workloads, hardware, and deployment contexts~\cite{kwon2023efficient,agrawal2024taming,lin2024towards,wang2016gunrock,shoeybi2019megatron,rasley2020deepspeed,pan2024survey,cao2023gpu}. Default policies often perform poorly: aggressive prefetching reduces stalls in LLM MoE inference but wastes bandwidth in graph analytics, while throughput-oriented scheduling improves vector search but worsens tail latency for interactive inference.

We are entering an age where AI agents increasingly optimize resource management policies. Finding optimal GPU resource management policies requires exploring a vast, workload-specific policy space, making manual tuning costly but suitable for AI-driven exploration. Agents can now generate complex code and optimize policy based on profiling results, such as compiler heuristics~\cite{alphaevolve2025}, datacenter scheduling, LLM serving~\cite{hamadanian2025glia,cheng2025barbarians}, training strategies~\cite{ding2025asap}, and CPU scheduling~\cite{dwivedula2025policysmith,zheng2025schedcp} rather than merely tuning parameters.

Effective GPU resource management requires programmable interfaces that are dynamic, safe, and full-stack. To leverage agent-driven exploration, interfaces need to be dynamic to support runtime policy updates without application restarts and service interrupt, programmable to express complex policies, and safe to enforce OS kernel isolation. To effectively manage resources across tenants, interfaces need to provide full-stack cross-application visibility and low-level hardware control, such as page fault handling and precise command scheduling.

Existing GPU resource management approaches fail to simultaneously meet these requirements. User-space frameworks~\cite{ng2023paella,shen2025xsched,gim2025pie,chen2025ktransformers} are safe and programmable but lack full-stack visibility and low-level hardware control. Kernel-driver modifications~\cite{kato2011timegraph,kato2012gdev,kernelintercept,fan2025gpreempt,wang2024gcaps} provide full-stack control but sacrifice safety and dynamic updates. Fixed configuration knobs exposed by existing APIs are safe but are not programmable.

We propose resolving this tradeoff by creating a programmable OS subsystem interface for GPU resource management that is full-stack, safe and dynamic. Our subsystem integrates directly within the GPU's OS kernel driver, which naturally controls cross-tenant hardware resources and coordinates across applications. To ensure safety and enable dynamic policy updates, we leverage eBPF, an OS kernel extensibility mechanism previously used in CPU scheduling~\cite{schedext-docs}, page caching~\cite{zussman2025cache_ext}, network management~\cite{bachl2021flow,Zhong22}, and storage systems~\cite{jia2023programmable}. eBPF allows verified user-defined callbacks to safely and dynamically replace hardcoded OS policies without service interruptions.

Alas, applying eBPF to GPU resource management is not straightforward. The GPU is a discrete processor with its own execution model and resources, and policy decisions made in the host take effect on a separate device. In particular, we identify three challenges: first, cross-device operations like migrating data or dispatching commands take milliseconds, too slow for synchronous callbacks and too costly for purely reactive decisions. Second, unlike independent subsystems on the CPU, limited device memory and bandwidth further couple memory and scheduling, as decisions about data placement directly affect which workloads can execute. Third, policy-relevant state is fragmented across the stack: GPU-internal execution details are opaque to the host, and application-level intent is hidden from the driver since user-space runtimes bypass the OS kernel.

To address these challenges, the key insight of \sys{} is to decouple policy triggers from their effects (\S\ref{sec:design-interfaces}): \sys{} models GPU resource management as programmable, asynchronous transitions on resource state machines. State machines define valid states and transitions: GPU memory chunks transition between host memory and device memory through eviction and prefetch; compute contexts transition through states such as running, ready via scheduling and preemption operations.
Unlike sched\_ext\cite{schedext-docs}, cache\_ext\cite{zussman2025cache_ext}, and XDP~\cite{bachl2021flow}, where each callback makes one synchronous decision at one fixed transition point, \sys{} policies attach to user-defined points across the stack, and each trigger can initiate multiple state transitions asynchronously. To ensure safety, policies only influence which transitions occur and when, but cannot create invalid transitions or violate resource invariants. \sys{} extends this model to GPU-internal execution, where a SIMT-aware verifier enables verified eBPF to run on GPU hardware for observation and policy execution (\S\ref{sec:verification}).

We implement \sys{} on Linux by extending the NVIDIA open GPU OS kernel modules~\cite{nvidia-open-gpu} and providing an eBPF runtime for both host driver and GPU device. We evaluate \sys{} across LLM inference, GNN training, vector search, and multi-tenant scenarios. \sys-based policies improve throughput by up to $4.8\times$ over framework offloading and $1.76\times$ over application-tuned baselines, and reduce tail latency by up to $2\times$, with low overhead and no application modifications or restarts.
In our experiments, AI agents generated all 59 evaluated eBPF policies, including eviction ordering, prefetch strategies, scheduling logic, and device-side work-stealing, with zero OS kernel panics, validating that eBPF's safety guarantees extend to the GPU driver and device. In summary, our contributions are:

\begin{itemize}[leftmargin=*,itemsep=0pt]
  \item We design a GPU-specific OS policy interface that models GPU resources as state machines whose asynchronous state transitions are directed by verified eBPF programs across the host-device boundary.
  \item We implement \sys, an eBPF runtime spanning host driver and GPU device, with a SIMT-aware verifier for device-side eBPF execution.
  \item We evaluate across four GPU workloads, achieving up to $4.8\times$ throughput improvement and $2\times$ tail latency reduction, with all 59 policies generated by AI agents and no OS kernel panics.
\end{itemize}
